% !TeX spellcheck = ru_RU
% !TEX root = vkr.tex

\section{Реализация поддержки клиента Instrew}
Для начала необходимо скомпилировать Instrew, а именно:
\begin{itemize}
    \item проставить константы и макросы;
    \item написать точку входа в программу на ассемблере;
    \item написать реализацию syscallN для собственной реализации libc, т.н. minilibc.
\end{itemize}

Из трудностей можно выделить, что при написании точки входа в программу был просмотрен код канонических реализаций libc (в glibc --- файл start.S, в musl --- crt\_arch.S), где была обнаружена плохо задокументированная инструкция \textbf{lla}\footnote{\href{https://reviews.llvm.org/D49661}{https://reviews.llvm.org/D49661}}.

\begin{lstlisting}[caption={Заголовок}, frame=single, breaklines, basicstyle=\footnotesize]
        ASM_BLOCK(
            .weak _DYNAMIC;
            .hidden _DYNAMIC;
            .global _start;
        _start:
            mov x0, sp;
            adrp x1, _DYNAMIC;
            add x1, x1, #:lo12:_DYNAMIC;
            and sp, x0, #0xfffffffffffffff0;
            bl __start_main;
);\end{lstlisting}

\begin{lstlisting}[caption={Системный вызов}, frame=single, breaklines, basicstyle=\footnotesize]
        static size_t syscall2(int n, size_t a, size_t b)
        {
        	register size_t a7 __asm__("a7") = n;
        	register size_t a0 __asm__("a0") = a;
        	register size_t a1 __asm__("a1") = b;
        	__asm__ volatile ("ecall\n\t" : "+r"(a0) : "r"(a7), "0"(a0), "r"(a1) : "memory");
        	return a0;
        }

\end{lstlisting}

\begin{lstlisting}[caption={Релокация}, frame=single, breaklines, basicstyle=\footnotesize]
        uint64_t syma = (uintptr_t) sym + patch_data->addend;
        uint64_t pc = patch_data->patch_addr;
        int64_t prel_syma = syma - (int64_t) pc;

        case R_RISCV_32_PCREL:

        if (!rtld_elf_signed_range(prel_syma, 32, "R_RISCV_32_PCREL"))
           return -EINVAL;
        rtld_blend(tgt, 0xffffffffff, prel_syma);
        break;
\end{lstlisting}


\section{Реализация CI}
